\section{Conclusiones y recomendaciones}
\label{sec:conclusiones}

Este trabajo propuso un modelo de dos niveles para la reforestación de 30 polígonos del matorral desértico de San Luis Potosí. El primero decide cuántas plantas de cada especie comprar en cada polígono para alcanzar la meta de 658 plantas por hectárea con una probabilidad prefijada, al menor costo y dentro del presupuesto. El segundo acomoda esas plantas en el tresbolillo para minimizar la competencia entre vecinas. Ambos se resolvieron con un método aproximado validado contra uno exacto, y se implementaron en una herramienta reproducible que CONAFOR puede usar con sus propios datos.

\paragraph{Compra} La restricción de probabilidad se reduce al cuantil de una binomial negativa, que la simulación Monte Carlo estima con un error del orden de \$1/ha frente al cálculo exacto. El resultado principal es que el presupuesto de \$12{,}500/ha no alcanza en ningún polígono: cumplir la meta con 95\,\% de confianza cuesta en promedio \$14{,}496/ha, 16\,\% más, es decir, \$365~mil adicionales para los 183.01~ha. La infactibilidad es estructural: el solo costo de cubrir el faltante, corregido por mortalidad, ya excede el presupuesto. Aceptar más riesgo apenas ayuda, pues pasar de 1\,\% a 50\,\% de probabilidad de incumplimiento ahorra 3.7\,\%. El margen de ajuste por especie, en cambio, es la principal palanca de costo: la compra óptima lo agota hacia los nopales, las especies con menor costo por planta sobreviviente. La compra resultante sobreplanta 9.8\,\% sobre la meta para absorber la mortalidad del primer año.

\paragraph{Distribución} Al redondear la compra a un mosaico de 52 posiciones, los 30 polígonos se reducen a seis claves de plantación. El algoritmo memético reduce la competencia entre 14.5\,\% y 16.3\,\% respecto al azar, con una variación casi nula entre corridas. El modelo entero no encontró una clave mejor en ningún caso y acota su distancia al óptimo en 4.5\,\% como máximo, una cota que refleja sobre todo la debilidad de la relajación lineal. Las claves eliminan casi toda vecindad entre plantas de la misma especie y usan al mezquite como separador y nodriza de las suculentas. Frente al orden vigente, la ganancia más grande no viene solo de optimizar, sino de dar el arreglo en dos dimensiones: la secuencia actual de 22 posiciones, repetida fila tras fila sin desfase, coloca a cada planta junto a dos de su especie y compite más que una plantación al azar. Las claves son robustas a los dos parámetros supuestos de la matriz de competencia, pero un error de $\pm$30\,\% en los rasgos de las especies las deja en promedio 6.6\,\% por encima de las que se obtendrían con los rasgos correctos.

\paragraph{Simulación de la siembra} En 1{,}000 siembras simuladas de una hectárea, el valor esperado de las plantas existentes es de 128.6 por hectárea, y para cumplir la meta de cada especie con 95\,\% de probabilidad hay que suministrar 651 plantas, con las que la hectárea queda en promedio con 722.4 plantas vivas al año. Cumplir cada especie con 95\,\% de probabilidad no equivale a cumplirlas todas: las diez a la vez se cumplen en 68.4\,\% de las réplicas. Con la clave 1, el valor esperado de la competencia al año es de 802.7 por hectárea, 18.0\,\% menos que una siembra al azar y 35.8\,\% menos que la secuencia vigente con filas alineadas. La matriz de transición de la clave tiene la diagonal nula, igual que la vigente; la diferencia entre ambas está en las vecinas de las filas contiguas, que solo una clave en dos dimensiones controla.

\paragraph{Recomendaciones} Para CONAFOR se desprenden cinco recomendaciones: (i) ajustar el presupuesto a alrededor de \$14{,}500/ha, o reducir la meta o la superficie a lo que el presupuesto actual cubre; (ii) aprovechar el margen de ajuste por especie hacia los nopales; (iii) entregar a las cuadrillas tablas de claves en dos dimensiones en lugar de una secuencia por fila; (iv) si se requiere que todas las especies cumplan a la vez con 95\,\% de probabilidad, repartir el riesgo entre ellas con $\alpha=0.005$ por especie, que por la desigualdad de Boole garantiza al menos 95\,\% de cumplimiento conjunto y eleva el presupuesto mínimo de \$14{,}496 a \$14{,}728/ha, apenas 1.6\,\%; y (v) continuar el monitoreo de supervivencia por especie y registrar la posición de las plantas, que son los datos que más limitan al modelo.

\paragraph{Trabajo futuro} El modelo puede extenderse en cuatro direcciones. La primera es un horizonte de varios años, con la cadena de Markov por posición de la \cref{subsec:markov}, ya que el monitoreo muestra mortalidad también en el segundo año. La segunda es la supervivencia correlacionada entre plantas por años secos y la incertidumbre en los conteos de existentes, ambas incorporables sorteándolas por escenario. La tercera es calibrar la matriz de competencia con datos de campo: con las coordenadas de cada planta, las funciones de correlación de pares estimarían la interacción entre especies que aquí se construyó con rasgos. La cuarta es acoplar los dos niveles, porque la competencia que reduce la distribución espacial debería elevar la supervivencia y, con ella, abaratar la compra.
